%% toc: P3 39 39.10
\chapter{Computable Presentations of Countable Structures}
\label{sec:p03-w03-p03-063-computable-presentations}

An abstract countable structure is not yet an object that an algorithm can
read.  To ask whether a structure is computable, one first fixes a concrete
copy whose elements are natural numbers.  One then asks whether its atomic
facts can be decided.  Two copies of the same structure can behave quite
differently, and some countable structures have no decidable copy at all.
This chapter fixes the conventions on which the short entry in
\cref{sec:p03-computable-model-theory} and the index-level result in
\cref{thm:p04-rational-order-uniform-index} depend.  These are the
signature hypothesis, the numerical coding of atomic and negated atomic
sentences, the domain convention including finite structures, indices and
computable isomorphisms.  It proves the facts these conventions require.

\paragraph{Inputs and conclusions.}
The input is a signature \(\Lang\) and a countable \(\Lang\)-structure, or a
concrete presentation of one.  The conclusions are of three kinds: whether a
given presentation is computable; whether the structure has some computable
presentation; and whether two computable presentations are computably
isomorphic.  The model-theoretic prerequisites are atomic diagrams and
parameter expansions (\cref{def:parameter-expansion,def:diagrams}),
isomorphisms (\cref{sec:maps-elementarity-01}) and, for the rational order,
classical back-and-forth (\cref{method:back-and-forth,thm:back-and-forth-test,thm:p03-rn-dlo-uniqueness}).
The computability prerequisites are the numerical and tuple coding
(\cref{sec:turing-machines-01}), partial
computable and computable functions (\cref{def:partial-computable}),
decidable and computably enumerable sets and relations
(\cref{sec:turing-machines-03}), the fixed
program numbering and its indices (\cref{def:p04-program-serialization}), and the
self-application set \(K\)
(\cref{sec:universal-reductions-03}, with \cref{thm:halting-undecidable}).

\paragraph{Selected treatment.}
The definitions follow the opening of Montalb\'an's \emph{Computable Structure
Theory: Within the Arithmetic}, in the author's draft of Part~1 compiled on
25 June 2019 \cite{Montalban2019CSTDraft}.  The book's version of record
\cite{Montalban2021CST} has not been consulted, and every locator below refers
to the 2019 draft.  Specifically: the computable-signature hypothesis follows
\cite[p.~xviii]{Montalban2019CSTDraft}; \(\omega\)-presentations, copies and
computability of a presentation follow Definition~I.1
\cite[pp.~1--2]{Montalban2019CSTDraft}; the atomic diagram as a coded object
follows Definition~I.2 \cite[p.~2]{Montalban2019CSTDraft}; presentations with
domain a subset of \(\N\), and the finite case, follow Sections~I.1.3 and
I.1.5 with Observation~I.3 \cite[pp.~3--5]{Montalban2019CSTDraft}; and
computable categoricity follows Definition~VIII.1
\cite[p.~123]{Montalban2019CSTDraft}.  The draft fixes only \emph{some}
effective enumeration of atomic formulas.  The explicit coding in
\cref{def:p03-w03-p03-063-sentence-coding} is the handbook's own instance of
it, and so is the identification of coded formulas with sentences of the
parameter expansion.  The draft asserts the equivalence of its two
computability conditions without proof \cite[p.~2]{Montalban2019CSTDraft};
\cref{thm:p03-w03-p03-063-equivalence} proves the corresponding statement.
The broader monograph of Ash and Knight \cite{AshKnight2000} is a separate
treatment, and nothing here is attributed to it.

\section{Signatures and the coding of literals}
\label{sec:p03-w03-p03-063-coding}

Throughout, signatures are single-sorted and first-order with equality, and
structures have nonempty domains, as in \cref{ch:maps-elementarity}.

\begin{hbconvention}{def:p03-w03-p03-063-computable-signature}{Computable signature}
Following \cite[p.~xviii]{Montalban2019CSTDraft}, a signature \(\Lang\) is
presented as three indexed families of symbols
\[
 \{\mathsf R_i:i\in I_R\},\qquad
 \{\mathsf f_i:i\in I_F\},\qquad
 \{\mathsf c_i:i\in I_C\},
\]
for relation, function and constant symbols, together with arity maps
\(a_R:I_R\to\N\) and \(a_F:I_F\to\N\).  Each of \(I_R,I_F,I_C\) is an
initial segment of \(\N\), finite or all of \(\N\).  The signature is
\emph{computable} when \(a_R\) and \(a_F\) are computable functions in the
sense of \cref{def:partial-computable}.  Every finite signature is
computable, since a map with finite domain is given by a finite table.  The
sans-serif \(\mathsf c_i\) are symbols of \(\Lang\).  They are distinct from
the parameter constants \(c_a\) of \cref{def:parameter-expansion}, which
name elements of a domain.
\index{computable signature}
\end{hbconvention}

The computability notions below are defined only for computable signatures.
For a noncomputable arity map, even the set of well-formed atomic formulas
can fail to be decidable.

\begin{hbconvention}{def:p03-w03-p03-063-sentence-coding}{Codes of terms and literals}
Fix a computable signature \(\Lang\) and variables \(x_0,x_1,\ldots\).
Write \(\langle a,b\rangle_{\N}\) for \(\operatorname{pair}_{\N}(a,b)\) and
\(\operatorname{code}_k\) for the tuple code of
\cref{sec:turing-machines-01}.  Codes of terms are given by
\begin{align*}
 \ulcorner x_j\urcorner&=\langle 0,j\rangle_{\N},\qquad
 \ulcorner \mathsf c_i\urcorner=\langle 1,i\rangle_{\N},\\
 \ulcorner \mathsf f_i(t_1,\ldots,t_k)\urcorner
 &=\bigl\langle 2,\langle i,\operatorname{code}_k(\ulcorner t_1\urcorner,\ldots,\ulcorner t_k\urcorner)\rangle_{\N}\bigr\rangle_{\N}
 \quad(k=a_F(i)).
\end{align*}
Codes of atomic formulas and their negations are given by
\begin{align*}
 \ulcorner t=s\urcorner&=\bigl\langle 3,\langle\ulcorner t\urcorner,\ulcorner s\urcorner\rangle_{\N}\bigr\rangle_{\N},\\
 \ulcorner \mathsf R_i(t_1,\ldots,t_k)\urcorner
 &=\bigl\langle 4,\langle i,\operatorname{code}_k(\ulcorner t_1\urcorner,\ldots,\ulcorner t_k\urcorner)\rangle_{\N}\bigr\rangle_{\N}
 \quad(k=a_R(i)),\\
 \ulcorner\neg\varphi\urcorner&=\langle 5,\ulcorner\varphi\urcorner\rangle_{\N}
 \quad(\varphi\text{ atomic}).
\end{align*}
An atomic formula or the negation of one is a \emph{literal}.

\emph{Sentences of a parameter expansion.}  Let \(D\subseteq\N\).  The
substitution \(x_j\mapsto c_j\) is a bijection from the literals of \(\Lang\)
all of whose variable indices lie in \(D\) onto the literal
\(\Lang(D)\)-sentences of \cref{def:parameter-expansion}.  Its inverse
replaces each parameter constant \(c_j\) by \(x_j\).  The code of a literal
\(\Lang(D)\)-sentence \(\sigma\) is the code of the literal it comes from:
\[
 \ulcorner\varphi(c_{j_1},\ldots,c_{j_m})\urcorner
 =\ulcorner\varphi(x_{j_1},\ldots,x_{j_m})\urcorner .
\]
This is the fixed serialization of atomic and negated atomic sentences used
for indices in \cref{sec:p04-busy-beaver-structures}.
\index{coding!of literals}
\end{hbconvention}

\begin{hblemma}{thm:p03-w03-p03-063-parsing}{proved=annex:p03-w03-p03-063-parsing}{Decidable syntax}
For a computable signature \(\Lang\), the sets \(\mathrm{Tm}\),
\(\mathrm{At}\) and \(\mathrm{Lit}\) of codes of terms, atomic formulas and
literals are decidable subsets of \(\N\).  Each code determines its
expression uniquely.  The relation ``\(x_j\) occurs in the expression with
code \(n\)'' is decidable on \(\mathrm{Tm}\cup\mathrm{Lit}\).  Every code of
an immediate subexpression is strictly smaller than the code of the
expression containing it.
\end{hblemma}

\section{Presentations and their atomic diagrams}
\label{sec:p03-w03-p03-063-presentations}

\begin{hbdefinition}{def:p03-w03-p03-063-presentations}{Presentations, coded diagrams and computability}
Let \(\Lang\) be a computable signature.
\begin{enumerate}
\item A \emph{\((\subseteq\omega)\)-presentation} is an \(\Lang\)-structure
  \(\mathcal M\) whose domain \(M\) is a nonempty subset of \(\N\).  It is
  an \emph{\(\omega\)-presentation} when \(M=\N\)
  \cite[Definition~I.1, p.~1]{Montalban2019CSTDraft}.  It is a
  \emph{\((\sqsubseteq\omega)\)-presentation} when \(M\) is an initial
  segment of \(\N\), that is, \(M=\{0,\ldots,k-1\}\) with \(k\geq1\) or
  \(M=\N\) \cite[pp.~3--4]{Montalban2019CSTDraft}.
\item A \emph{presentation of}, or \emph{copy of}, an \(\Lang\)-structure
  \(\mathcal A\) is a \((\subseteq\omega)\)-presentation isomorphic to
  \(\mathcal A\).
\item The \emph{coded atomic diagram} of \(\mathcal M\) is
  \[
   \ulcorner\operatorname{Diag}(\mathcal M)\urcorner
   =\{\ulcorner\sigma\urcorner:\sigma\in\operatorname{Diag}(\mathcal M)\}
   \subseteq\N,
  \]
  where \(\operatorname{Diag}(\mathcal M)\) is the set of literal
  \(\Lang(M)\)-sentences true in \(\mathcal M_M\)
  (\cref{def:diagrams}) and codes are those of
  \cref{def:p03-w03-p03-063-sentence-coding}.
\item \(\mathcal M\) is \emph{computable} when
  \(\ulcorner\operatorname{Diag}(\mathcal M)\urcorner\) is decidable
  (\cref{sec:turing-machines-03}).  A structure \emph{has a computable
  copy} when some copy of it is computable.
\item The interpretations of \(\mathcal M\) are \emph{uniformly computable}
  when there are computable total functions
  \(\rho,\zeta:\N^2\to\N\) and \(\gamma:\N\to\N\) such that, for
  \(\vec n\in M^k\),
  \begin{gather*}
   \rho(i,\operatorname{code}_k(\vec n))=1\iff \vec n\in\mathsf R_i^{\mathcal M}
   \quad(i\in I_R,\ k=a_R(i)),\\
   \zeta(i,\operatorname{code}_k(\vec n))=\mathsf f_i^{\mathcal M}(\vec n)
   \quad(i\in I_F,\ k=a_F(i)),
  \end{gather*}
  and \(\gamma(i)=\mathsf c_i^{\mathcal M}\) for \(i\in I_C\).
  This is the handbook form of the condition that the set
  \(\tau^{\mathcal M}\) of \cite[Definition~I.1, p.~2]{Montalban2019CSTDraft}
  be computable.
\end{enumerate}
A code of a literal that mentions a variable \(x_j\) with \(j\notin M\) is
never in \(\ulcorner\operatorname{Diag}(\mathcal M)\urcorner\).  This
matches the draft's convention for domains that are proper subsets of \(\N\)
\cite[Section~I.1.3, p.~3]{Montalban2019CSTDraft}.
\index{presentation!omega-presentation}
\index{atomic diagram!coded}
\end{hbdefinition}

For an \(\omega\)-presentation, this is the notion of
\cref{def:p03-computable-presentation}, with the atomic diagram read through
the coding above.  A finite structure has no \(\omega\)-presentation, since
\(\N\) is infinite, but it has \((\sqsubseteq\omega)\)-presentations.
\Cref{thm:p03-w03-p03-063-domain} describes exactly what the choice of
domain changes.

\begin{hbtheorem}{thm:p03-w03-p03-063-equivalence}{proved=annex:p03-w03-p03-063-equivalence}{Three forms of a computable presentation}
Let \(\mathcal M\) be a \((\subseteq\omega)\)-presentation in a computable
signature.  The following are equivalent.
\begin{enumerate}
\item \(\ulcorner\operatorname{Diag}(\mathcal M)\urcorner\) is decidable.
\item The set of codes of true atomic \(\Lang(M)\)-sentences is decidable,
  and so is \(M\).
\item \(M\) is decidable and the interpretations of \(\mathcal M\) are
  uniformly computable.
\end{enumerate}
Moreover, in each case \(M\) is decidable.  Let \(a_0<a_1<\cdots\) list
\(\mathrm{At}\).  The binary sequence \(D(\mathcal M)\) with
\(D(\mathcal M)(i)=1\) iff \(a_i\) codes a sentence in
\(\operatorname{Diag}(\mathcal M)\) is computable iff the three conditions
hold.  This is the atomic diagram \(D(\mathcal M)\in2^{\N}\) of
\cite[Definition~I.2, p.~2]{Montalban2019CSTDraft}, for the enumeration
\((a_i)\).
\end{hbtheorem}

\begin{hbmethod}{def:p03-w03-p03-063-decide-literal}{Deciding a literal}
\emph{Input:} the decision procedure for \(M\), the functions
\(\rho,\zeta,\gamma\) of a computable presentation, and a number \(n\).
\emph{Output:} whether \(n\in\ulcorner\operatorname{Diag}(\mathcal M)\urcorner\).
\begin{enumerate}
\item Decide whether \(n\in\mathrm{Lit}\); if not, answer no.
\item List the indices \(j\) with \(x_j\) occurring in \(n\); if one lies
  outside \(M\), answer no.
\item Evaluate every term by recursion on codes:
  \(x_j\mapsto j\), \(\mathsf c_i\mapsto\gamma(i)\) and
  \(\mathsf f_i(\vec t)\mapsto\zeta(i,\operatorname{code}_k(\text{values of }\vec t))\).
\item Answer according to equality of values, or according to \(\rho\),
  and reverse the answer for a negated atomic code.
\end{enumerate}
\emph{Side conditions:} the signature must be computable, so that steps~1
and~3 can read arities; and \(M\) must be decidable, so that step~2
terminates.  The procedure decides a fixed presentation.  It says nothing
about other copies of the same structure.
\end{hbmethod}

\section{Indices and computable isomorphisms}
\label{sec:p03-w03-p03-063-indices}

\begin{hbdefinition}{def:p03-w03-p03-063-index}{Indices and computable isomorphisms}
Let \(\mathcal M\) be a computable \((\subseteq\omega)\)-presentation.  An
\emph{index} of \(\mathcal M\) is a number \(e\) such that \(\varphi_e\), in
the numbering of \cref{def:p04-program-serialization}, is total, takes only the
values \(0\) and \(1\), and is the characteristic function of
\(\ulcorner\operatorname{Diag}(\mathcal M)\urcorner\).  An isomorphism
\(g:\mathcal A\to\mathcal B\) between \((\subseteq\omega)\)-presentations is
\emph{computable} when some partial computable function is defined on all of
\(A\) and agrees with \(g\) there.  Two presentations are \emph{computably
isomorphic} when a computable isomorphism exists between them.
\index{computable presentation!index}
\index{computable isomorphism}
\end{hbdefinition}

An index depends on the program numbering and on the coding of
\cref{def:p03-w03-p03-063-sentence-coding}.  It is a name of one
presentation, not of an isomorphism type.  By
\cref{thm:p03-w03-p03-063-equivalence}, the presentation can be recovered
from its index: the domain, the relations, the functions and the constants
are all read off the decided diagram.  The draft defines the \emph{index set}
of a class of structures as the set of indices of computable functions that
are diagrams of \(\omega\)-presentations of members of the class
\cite[p.~124, note]{Montalban2019CSTDraft}.

\begin{hblemma}{thm:p03-w03-p03-063-pullback}{proved=annex:p03-w03-p03-063-pullback}{Transport along computable bijections}
Let \(\mathcal M\) be a computable \(\omega\)-presentation and
\(g:\N\to\N\) a computable bijection.  Then \(g^{-1}\) is computable.  The
\emph{pullback} \(g^{-1}(\mathcal M)\), with domain \(\N\) and
\begin{gather*}
 \mathsf R_i^{g^{-1}(\mathcal M)}=\{\vec n:g\vec n\in\mathsf R_i^{\mathcal M}\},\\
 \mathsf c_i^{g^{-1}(\mathcal M)}=g^{-1}(\mathsf c_i^{\mathcal M}),\\
 \mathsf f_i^{g^{-1}(\mathcal M)}(\vec n)=g^{-1}\bigl(\mathsf f_i^{\mathcal M}(g\vec n)\bigr),
\end{gather*}
is a computable \(\omega\)-presentation, and \(g\) is an isomorphism
\(g^{-1}(\mathcal M)\to\mathcal M\).  Consequently an
\(\omega\)-presentation computably isomorphic to a computable one is itself
computable.  Composites and inverses of computable isomorphisms between
\(\omega\)-presentations are computable.
\end{hblemma}

By the last clause, \cref{def:p03-computably-categorical} agrees with
\cite[Definition~VIII.1, p.~123]{Montalban2019CSTDraft}.  The handbook
definition asks for a computable isomorphism from a fixed computable
\(\mathcal A\) to each computable copy.  The draft asks for one between any
two computable copies.  Composing through \(\mathcal A\) and inverting turns
either requirement into the other.

\section{Domains and computable copies}
\label{sec:p03-w03-p03-063-domains}

\begin{hbtheorem}{thm:p03-w03-p03-063-domain}{proved=annex:p03-w03-p03-063-domain}{What the domain convention changes}
Let \(\Lang\) be a computable signature.
\begin{enumerate}
\item If \(\mathcal M\) is a computable \((\subseteq\omega)\)-presentation
  with infinite domain \(M\), list \(M=\{m_0<m_1<\cdots\}\) and put
  \(h(i)=m_i\).  Then the \(\omega\)-presentation \(\mathcal A\) obtained by
  transporting \(\mathcal M\) along \(h\) is computable, and
  \(h:\mathcal A\to\mathcal M\) is a computable isomorphism whose inverse
  \(M\to\N\) is also computable \cite[Observation~I.3, pp.~3--4]{Montalban2019CSTDraft}.
\item If \(\Lang\) is finite, every finite \(\Lang\)-structure with \(k\)
  elements has a computable presentation with domain \(\{0,\ldots,k-1\}\);
  this is a \((\sqsubseteq\omega)\)-presentation.
\item There is a computable infinite signature, consisting of unary relation
  symbols \(\mathsf R_i\) for \(i\in\N\), and a one-element structure with no
  computable \((\subseteq\omega)\)-presentation.
\end{enumerate}
\end{hbtheorem}

Clause~(a) is why the entry in \cref{sec:p03-computable-model-theory} can use
domain \(\N\) for infinite structures without loss.  Clause~(b) is the
correct statement for finite structures, which have no
\(\omega\)-presentation.  Clause~(c) shows that finiteness of the domain alone
does not make a structure computable: in an infinite signature, a single
element can carry a noncomputable set of atomic facts.

\begin{hbtheorem}{thm:p03-w03-p03-063-no-copy}{proved=annex:p03-w03-p03-063-no-copy}{An abstract countable structure need not have a computable copy}
Let \(\Lang=\{E\}\) with \(E\) binary.  For an infinite \(X\subseteq\N\), let
\(\mathcal C_X\) be the graph that is a disjoint union of cycles, with
exactly one cycle of length \(n+3\) for each \(n\in X\) and no other
vertices.  If \(\mathcal C_X\) has a computable copy, then \(X\) is
computably enumerable.  In particular, \(\mathcal C_{\N\setminus K}\) is a
countably infinite \(\Lang\)-structure with no computable copy.
\end{hbtheorem}

The structure \(\mathcal C_{\N\setminus K}\) exists, is countable, and has
\(\omega\)-presentations, but none of them is decidable.  The existence of an
abstract countable model is therefore never by itself a computable
presentation.  A claim that a structure is computable must exhibit, or prove
the existence of, a presentation together with a decision procedure for its
coded diagram.

\section{Copies of one structure that differ computationally}
\label{sec:p03-w03-p03-063-copies}

\begin{hbtheorem}{thm:p03-w03-p03-063-omega-copies}{proved=annex:p03-w03-p03-063-omega-copies}{Computable copies of \((\N,<)\) that are not computably isomorphic}
Let \(\Lang=\{<\}\).
\begin{enumerate}
\item There is a computable \(\omega\)-presentation \(\mathcal M=(\N,<_M)\)
  isomorphic to \((\N,<)\) such that the \(<_M\)-immediate successor of
  \(2n\) is \(2n+2\) exactly when \(n\notin K\).  Hence \(\mathcal M\) and
  \((\N,<)\) are not computably isomorphic, and \((\N,<)\) is not computably
  categorical.
\item There is an \(\omega\)-presentation of \((\N,<)\) that is not
  computable.
\end{enumerate}
\end{hbtheorem}

Clause~(a) follows the example of \cite[Section~I.1.2, pp.~2--3]{Montalban2019CSTDraft}.
The draft orders by \(\leq\) and codes its halting set \(0'\).  Here the
strict order and the self-application set \(K\) of
\cref{sec:universal-reductions-03} are used, so that the proof runs
entirely on Part~IV results.  Clause~(b) is a separate, elementary
observation.  It shows that computability is a property of a presentation,
not of an isomorphism type.  The contrast with
\cref{thm:p03-q-computably-categorical} is the point.  For the dense order,
the classical back-and-forth argument can be run effectively, and uniformly
in indices
(\cref{thm:p04-rational-order-uniform-index}).  For \((\N,<)\), the
immediate-successor function is computable in the standard copy but not in
the copy \(\mathcal M\), and a computable isomorphism would transfer it, so no
such effective construction is possible.

\begin{hbexample}{def:p03-w03-p03-063-worked}{Coding and deciding one literal}
Take \(\Lang=\{\mathsf R_0\}\) with \(a_R(0)=2\), \(I_F=I_C=\varnothing\),
and write \(\mathsf R_0\) as \(<\).  In the standard presentation
\((\N,<)\), consider the \(\Lang(\N)\)-sentence \(c_3<c_5\).  It comes from
the formula \(x_3<x_5\), and
\begin{gather*}
 \ulcorner x_3\urcorner=\langle0,3\rangle_{\N}=9,\qquad
 \ulcorner x_5\urcorner=\langle0,5\rangle_{\N}=20,\\
 \operatorname{code}_2(9,20)=\langle9,20\rangle_{\N}=455,
\end{gather*}
so
\[
 \ulcorner c_3<c_5\urcorner
 =\bigl\langle4,\langle0,455\rangle_{\N}\bigr\rangle_{\N}
 =\langle4,104195\rangle_{\N}=5428872095 .
\]
Since \(3<5\), this number lies in
\(\ulcorner\operatorname{Diag}(\N,<)\urcorner\), while
\(\ulcorner\neg(c_3<c_5)\urcorner=\langle5,5428872095\rangle_{\N}\) does not.
In the procedure of \cref{def:p03-w03-p03-063-decide-literal}, step~1
recognizes the tag \(4\), the relation index \(0\) and the arity \(2\).
Step~2 finds the variable indices \(3\) and \(5\) in \(M=\N\).  Step~3
evaluates the two terms to \(3\) and \(5\).  Step~4 applies
\(\rho(0,\operatorname{code}_2(3,5))=1\).  Every step is a bounded
computation.  Running the same procedure on the presentation \(\mathcal M\)
of \cref{thm:p03-w03-p03-063-omega-copies} gives a different diagram, because
there \(\rho\) compares the keys of the construction rather than the numbers
themselves.
\end{hbexample}

\begin{hbwarning}{def:p03-w03-p03-063-limits}{Limits of the notions}
\begin{enumerate}
\item The notions require a computable signature.  Without one, the coding
  of \cref{def:p03-w03-p03-063-sentence-coding} need not have a decidable
  set of well-formed codes.
\item Computability belongs to a presentation.  A computable structure has
  noncomputable copies (\cref{thm:p03-w03-p03-063-omega-copies}).
\item Having a computable copy is a property of the isomorphism type, and it
  can fail for countable structures (\cref{thm:p03-w03-p03-063-no-copy}).
\item A finite structure has no \(\omega\)-presentation.  In an infinite
  computable signature, it need not have any computable presentation
  (\cref{thm:p03-w03-p03-063-domain}).
\item An index names a presentation relative to fixed codings.  Changing the
  program numbering or the literal coding changes the indices, and it does not
  change which presentations are computable.
\item The draft reports that the index set of computably categorical
  structures is \(\Pi^1_1\)-complete, a result it attributes to Downey,
  Kach, Lempp, Lewis-Pye, Montalb\'an and Turetsky
  \cite[p.~124]{Montalban2019CSTDraft}.  The handbook neither proves nor uses
  this.  It is recorded only as the reason why computable categoricity
  admits no simple structural test in general.
\end{enumerate}
\end{hbwarning}

\section{Relationships and notation}
\label{sec:p03-w03-p03-063-relations}

\begin{table}[ht]
\centering
\small
\begin{tabular}{p{0.38\linewidth}p{0.52\linewidth}}
\toprule
Draft \cite{Montalban2019CSTDraft} & Handbook \\
\midrule
vocabulary \(\tau\), arities \(a_R,a_F\) (p.~xviii) & signature \(\Lang\), \cref{def:p03-w03-p03-063-computable-signature} \\
\(\omega\)-presentation, copy (Def.~I.1) & presentation with domain \(\N\), \cref{def:p03-w03-p03-063-presentations} \\
\(\tau^{\mathcal M}\) computable (Def.~I.1) & uniformly computable interpretations \(\rho,\zeta,\gamma\) \\
\(\varphi^{\mathrm{at}}_i\), \(D(\mathcal M)\in2^{\N}\) (Def.~I.2) & \(\mathrm{At}\), \(\ulcorner\operatorname{Diag}(\mathcal M)\urcorner\subseteq\N\), \cref{thm:p03-w03-p03-063-equivalence} \\
\((\subseteq\omega)\)-, \((\sqsubseteq\omega)\)-presentations (pp.~3--4) & the same names, \cref{def:p03-w03-p03-063-presentations} \\
computably categorical (Def.~VIII.1) & \cref{def:p03-computably-categorical}, via \cref{thm:p03-w03-p03-063-pullback} \\
halting set \(0'\) (p.~3) & self-application set \(K\), \cref{sec:universal-reductions-03} \\
\bottomrule
\end{tabular}
\caption{Notation of the selected treatment and of the handbook}
\label{tab:p03-w03-p03-063-bridge}
\end{table}

The classical theorem that countable dense linear orders without endpoints
are isomorphic (\cref{thm:p03-rn-dlo-uniqueness}) is a statement about
isomorphism types.  Its effective form (\cref{thm:p03-q-computably-categorical})
is a statement about computable presentations, and it uses the decidability of
both coded diagrams at every stage of the back-and-forth construction.  The
draft isolates the abstract mechanism: a set of pairs of tuples with the
back-and-forth property yields an isomorphism computable from an enumeration
of that set \cite[Lemma~III.15, p.~40]{Montalban2019CSTDraft}.  Finite model
theory, the subject of the surrounding chapters, concerns finite structures
as inputs to algorithms, and \cref{thm:p03-w03-p03-063-domain} records how
that setting fits under the present conventions.
